% !TeX program = xelatex
% ============================================================
% rp/rp.tex — Руководство программиста
% ГОСТ 19.504-79 «Руководство программиста. Требования к
% содержанию и оформлению»
% ============================================================
\documentclass[12pt,a4paper]{extarticle}

% 1. Подключение общей преамбулы (пакеты)
\input{../common/preamble}

% 2. Определение переменных для конкретного документа
\newcommand{\docname}{Programmer's Manual}
\newcommand{\doccode}{\hseDocCodeBase\ 33 01-1}
\newcommand{\doccodelu}{\hseDocCodeBase\ 33 01-1-ЛУ}

% 3. Подключение общих команд (проект, студент, руководитель)
\input{../common/commands}

% 4. Подключение стилей (колонтитулы, заголовки)
\input{../common/styles}

\begin{document}
\sloppy

% 5. Генерация титульных листов по шаблону
\input{../common/title_template}


% ============================================================
% ANNOTATION
% ============================================================
\section*{ANNOTATION}
\addcontentsline{toc}{section}{ANNOTATION}

This document, \enquote{\hseProjectName. Programmer's Manual}, contains the information required for the development, configuration, launch, verification, and maintenance of the software complex. It describes the purpose and conditions of use of the program, the characteristics of its architecture, the procedure for addressing the program, the organization of input and output data, as well as the diagnostic messages generated by the system.

% Hint: specify who the manual is intended for (which services,
% libraries, and configurations its readers develop).
The manual is intended for programmers and engineers who develop the services, shared libraries, software contracts, and infrastructure configurations of the project.

The document has been developed in accordance with the requirements of GOST 19.504-79\cite{gost19504}.

\clearpage

% ============================================================
% LIST OF APPLICABLE GOST STANDARDS
% ============================================================
\noindent
This document has been developed in accordance with the requirements of:

\begin{enumerate}[label=\arabic*), leftmargin=1.25cm]
  \item GOST 19.101-77 Types of programs and program documents\cite{gost19101}.
  \item GOST 19.103-77 Designation of programs and program documents\cite{gost19103}.
  \item GOST 19.104-78 Basic inscriptions\cite{gost19104}.
  \item GOST 19.105-78 General requirements for program documents\cite{gost19105}.
  \item GOST 19.106-78 Requirements for program documents produced by
        printed means\cite{gost19106}.
  \item GOST 19.504-79 Programmer's manual. Requirements for content
        and presentation\cite{gost19504}.
\end{enumerate}

Amendments to this document are made in accordance with GOST\,19.604-78\cite{gost19604}.

\clearpage

% ============================================================
% TABLE OF CONTENTS
% ============================================================
\hypersetup{linkcolor=black}
\tableofcontents
\hypersetup{linkcolor=blue}
\thispagestyle{fancy}
\clearpage

% ============================================================
% 1. PURPOSE AND CONDITIONS OF USE OF THE PROGRAM
% ============================================================
\section{PURPOSE AND CONDITIONS OF USE OF THE PROGRAM}

\subsection{Purpose of the program}

% Hint: describe the purpose of the program and the range of tasks it solves.
\begin{hseExample}
The program \enquote{\projectname} is intended for \hseFill{specify the purpose of the program}. The system provides \hseFill{categories of users} with software interfaces for \hseFill{list of tasks to be solved}.

The program is used as \hseFill{the role of the program within the overall system} and as a basis for the further development of \hseFill{directions of functionality development}.
\end{hseExample}

\subsection{Functions performed by the program}

% Hint: list the main functions of the program (an itemized list).
\begin{hseExample}
The software complex performs the following main functions:

\begin{itemize}[leftmargin=1.25cm]
    \item \hseFill{the first main function of the program};
    \item \hseFill{the second main function of the program};
    \item \hseFill{the third main function of the program};
    \item collection of metrics, traces, and structured logs.
\end{itemize}
\end{hseExample}

\subsection{Conditions required for running the program}

% Hint: specify the environment in which the program is launched
% (containers, OS, development environment, etc.).
\begin{hseExample}
The program is designed to run in \hseFill{the program execution environment}. For local development, \hseFill{local development tools} are used.
\end{hseExample}

\subsubsection{Hardware requirements}

% Hint: minimum/recommended requirements for the processor, RAM,
% disk space, and network.
\begin{hseExample}
For running the program locally, a workstation with the following characteristics is recommended:

\begin{itemize}[leftmargin=1.25cm]
    \item processor: at least \hseFill{number of cores and architecture};
    \item random-access memory: at least \hseFill{amount of RAM};
    \item disk space: at least \hseFill{amount of disk space} of free space for the source code, dependencies, and data;
    \item network access to \hseFill{repository and external resources}.
\end{itemize}

For a production deployment, the components must be placed in \hseFill{the target deployment environment} with dedicated resources for \hseFill{the DBMS and infrastructure services}.
\end{hseExample}

\subsubsection{Software requirements}

% Hint: list the OS, runtimes, DBMS, build tools, and other software.
\begin{hseExample}
The following software is required to develop and run the program:

\begin{itemize}[leftmargin=1.25cm]
    \item operating system \hseFill{supported operating systems};
    \item Git for obtaining the source code;
    \item \hseFill{containerization or virtualization tools};
    \item \hseFill{programming language and its version};
    \item \hseFill{DBMS and infrastructure dependencies};
    \item \hseFill{additional development tools}.
\end{itemize}
\end{hseExample}

\subsubsection{Personnel requirements}

% Hint: specify the qualifications of the programmer (languages, technologies, skills).
\begin{hseExample}
The programmer who develops and maintains the program must have practical skills in working with \hseFill{languages, frameworks, and tools} and with the Git version control system. Modifying \hseFill{an adjacent layer of the system} additionally requires skills in working with \hseFill{additional technologies and tools}.
\end{hseExample}

\clearpage

% ============================================================
% 2. CHARACTERISTICS OF THE PROGRAM
% ============================================================
\section{CHARACTERISTICS OF THE PROGRAM}

\subsection{Description of the main characteristics of the program}

% Hint: describe the architecture, the approaches used, and the technologies.
\begin{hseExample}
The program \enquote{\projectname} is implemented as \hseFill{the type of software complex}. The server side is built on \hseFill{languages and frameworks} and uses \hseFill{the architectural style}. Interaction between components is performed via \hseFill{interaction protocols}.

The system uses the following architectural approaches:

\begin{itemize}[leftmargin=1.25cm]
    \item separation of components by areas of responsibility;
    \item \hseFill{the second architectural approach};
    \item \hseFill{the third architectural approach};
    \item data schema migrations via \hseFill{the migration tool}.
\end{itemize}
\end{hseExample}

\subsection{Operating mode of the program}

% Hint: specify the operating mode (continuous/batch/interactive) and
% the permitted launch scenarios.
\begin{hseExample}
In normal operation the program runs continuously. \hseFill{describe the operation of the components in normal mode}.

For development, the following launch modes are permitted:

\begin{itemize}[leftmargin=1.25cm]
    \item launching only the infrastructure dependencies;
    \item launching a selected component of the program;
    \item launching all components simultaneously;
    \item launching a test environment for integration checks.
\end{itemize}
\end{hseExample}

\subsection{Composition and purpose of the components}

% Hint: list the main components of the program and their purpose.
\begin{hseExample}
The composition and purpose of the main components of the program are given in Table~\ref{tab:rp_components}.

% The table is typeset via \captionof (not \begin{table}[H]): floating
% environments inside the yellow sample block are not allowed.
\begin{center}
\small
\captionof{table}{Main components of the program}
\label{tab:rp_components}
\begin{tabularx}{\textwidth}{|>{\raggedright\arraybackslash}p{0.34\textwidth}|X|}
\hline
\textbf{Component} & \textbf{Purpose} \\
\hline
\hseFill{component name} & \hseFill{component purpose}. \\
\hline
\hseFill{component name} & \hseFill{component purpose}. \\
\hline
\end{tabularx}
\end{center}
\end{hseExample}

\subsection{Means of controlling correct execution}

% Hint: describe the means of controlling correct operation (tests,
% linters, health checks, logs, metrics, traces).
\begin{hseExample}
The correctness of the program's operation is controlled by the following means:

\begin{itemize}[leftmargin=1.25cm]
    \item automated unit and integration tests via \hseFill{the test framework};
    \item formatting, static analysis, and type checks via \hseFill{linters and analysers};
    \item health checks of the program's components;
    \item metrics, traces, and structured logs;
    \item \hseFill{additional means of control}.
\end{itemize}
\end{hseExample}

\clearpage

% ============================================================
% 3. ADDRESSING THE PROGRAM
% ============================================================
\section{ADDRESSING THE PROGRAM}

\subsection{Obtaining the source code}

% Hint: specify the location of the repository and the cloning commands.
\begin{hseExample}
The source code of the program is located in the remote Git repository \hseFill{specify the Git hosting and the repository name}.

To obtain the source code, run:

\begin{verbatim}
git clone <repository URL>
cd <project directory>
\end{verbatim}
\end{hseExample}

\subsection{Preparing the development environment}

% Hint: describe the preparation of configurations and the installation of dependencies.
\begin{hseExample}
Before launching the program, install the system dependencies specified in subsection 1.3.2 and prepare the environment variable files. Example \texttt{.env.example} configurations are located in \hseFill{directories with example configurations}.

The minimum preparation procedure:

\begin{enumerate}[label=\arabic*), leftmargin=1.25cm]
    \item copy the example configurations into working \texttt{.env} files;
    \item install the dependencies with the command \hseFill{dependency installation command};
    \item \hseFill{additional environment preparation steps};
    \item apply the database migrations.
\end{enumerate}
\end{hseExample}

\subsection{Launching the program}

% Hint: provide the commands for launching the program and its components
% (infrastructure dependencies, an individual component, the whole environment).
\begin{hseExample}
To launch the infrastructure dependencies, use the command:

\begin{verbatim}
<infrastructure launch command>
\end{verbatim}

To launch the program and all of its components:

\begin{verbatim}
<program launch command>
\end{verbatim}

The local environment is stopped with the command:

\begin{verbatim}
<stop command>
\end{verbatim}
\end{hseExample}

\subsection{Configuration}

% Hint: describe the configuration parameters (environment variables,
% configuration files) and their default values.
\begin{hseExample}
The program parameters are set by environment variables and configuration files. The main parameters include connection strings to \hseFill{infrastructure dependencies}, observability parameters, service keys, and component launch flags.
\end{hseExample}

\subsection{Running checks}

% Hint: provide the commands for running tests and other checks.
\begin{hseExample}
To run tests, static checks, and formatting, the following commands are used:

\begin{verbatim}
<test run command>
<static check command>
<formatting command>
\end{verbatim}

For a smoke check of a locally running environment:

\begin{verbatim}
<smoke check command>
\end{verbatim}
\end{hseExample}

\clearpage

% ============================================================
% 4. INPUT AND OUTPUT DATA
% ============================================================
\section{INPUT AND OUTPUT DATA}

\subsection{Organization of input data}

% Hint: describe the sources, formats, and composition of the input data.
\begin{hseExample}
The input data of the program is received through \hseFill{input data interfaces and channels}, environment variables, configuration files, and migration data. The main input data formats are given in Table~\ref{tab:rp_input_data}.

% The table is typeset via \captionof (not \begin{table}[H]): floating
% environments inside the yellow sample block are not allowed.
\begin{center}
\small
\captionof{table}{Input data of the program}
\label{tab:rp_input_data}
\begin{tabularx}{\textwidth}{|>{\raggedright\arraybackslash}p{0.28\textwidth}|X|}
\hline
\textbf{Source} & \textbf{Data composition} \\
\hline
\hseFill{API interface} & \hseFill{composition of request data}. \\
\hline
Environment variables & Connection strings to \hseFill{infrastructure dependencies}, observability parameters, keys, and component launch flags. \\
\hline
Database migrations & \hseFill{migration tool} scripts for changing the database structure. \\
\hline
\end{tabularx}
\end{center}
\end{hseExample}

\subsection{Organization of output data}

% Hint: describe the channels, formats, and composition of the output data.
\begin{hseExample}
The output data is generated in the form of \hseFill{output data channels}, database records, metrics, traces, and logs. The main types of output information are given in Table~\ref{tab:rp_output_data}.

% The table is typeset via \captionof (not \begin{table}[H]): floating
% environments inside the yellow sample block are not allowed.
\begin{center}
\small
\captionof{table}{Output data of the program}
\label{tab:rp_output_data}
\begin{tabularx}{\textwidth}{|>{\raggedright\arraybackslash}p{0.28\textwidth}|X|}
\hline
\textbf{Channel} & \textbf{Data composition} \\
\hline
\hseFill{API channel} & \hseFill{composition of responses and errors}. \\
\hline
Database & Business data of the program, migration history, and service structures. \\
\hline
Observability & Metrics, traces, error events, and structured logs. \\
\hline
\end{tabularx}
\end{center}
\end{hseExample}

\subsection{Data constraints and validation}

% Hint: describe the levels of validation and the constraints on the data.
\begin{hseExample}
Validation of the input data is performed at several levels:

\begin{itemize}[leftmargin=1.25cm]
    \item \hseFill{validation schemas} check the types, mandatory presence, and format of the fields;
    \item the domain layer checks business invariants: \hseFill{examples of business rules};
    \item database constraints and indexes protect against violations of data integrity;
    \item \hseFill{additional data protection mechanisms}.
\end{itemize}
\end{hseExample}

\clearpage

% ============================================================
% 5. MESSAGES
% ============================================================
\section{MESSAGES}

% Hint: describe the types of messages generated and fill in the table
% of diagnostic messages (code/text, cause, programmer's actions).
\begin{hseExample}
The program generates messages of several types: \hseFill{types of responses and error codes}, log records, health-check messages, and monitoring events. The messages are intended for the programmer, the system operator, and the calling software clients. The main diagnostic messages are given in Table~\ref{tab:rp_messages}.
\end{hseExample}

% The longtable remains outside the yellow sample block: breakable tables inside
% hseExample are not allowed.
\begin{small}
\begin{longtable}{|>{\raggedright\arraybackslash}p{0.19\textwidth}|>{\raggedright\arraybackslash}p{0.21\textwidth}|>{\raggedright\arraybackslash}p{0.48\textwidth}|}
\caption{Main diagnostic messages and programmer's actions}
\label{tab:rp_messages}\\
\hline
\textbf{Message/code} & \textbf{Cause} & \textbf{Programmer's actions} \\
\hline
\endfirsthead
\hline
\textbf{Message/code} & \textbf{Cause} & \textbf{Programmer's actions} \\
\hline
\endhead
\hline
\endfoot
\texttt{400 Bad Request} & Incorrect format of the input data or request parameters. & Check the request schema, the client-side validation, and the boundary-value tests. \\
\hline
\texttt{404 Not Found} & The requested entity was not found. & Check the identifier, the data, and the correctness of the route. \\
\hline
\texttt{503 Service Unavailable} & A dependent service, database, or external component is unavailable. & Check the health checks, the environment variables, the network settings, and the logs of the corresponding service. \\
\hline
\texttt{Health check not ready} & The component has started but is not ready to accept traffic. & Check the readiness endpoint, the migrations, and the connections to the dependencies. \\
\hline
\hseFill{message or code} & \hseFill{cause} & \hseFill{programmer's actions} \\
\hline
\end{longtable}
\end{small}

\clearpage
% The pages of terms, references, and appendices are printed without the bottom
% stamp (as in real document sets) — only the sheet number and the document code at the top.
\pagestyle{appendix}

% ============================================================
% LIST OF TERMS AND DEFINITIONS
% ============================================================
\section*{LIST OF TERMS AND DEFINITIONS}
\addcontentsline{toc}{section}{LIST OF TERMS AND DEFINITIONS}

% Hint: list the terms and abbreviations used in the document
% (in alphabetical order).
\begin{enumerate}[label=\arabic*., leftmargin=1.25cm]
  \item \textbf{\hseFill{term}} --- \hseFill{definition of the term}.
\end{enumerate}

\clearpage

% ============================================================
% LIST OF REFERENCES
% ============================================================
\section*{LIST OF REFERENCES}
\addcontentsline{toc}{section}{LIST OF REFERENCES}

% Hint: before the GOST standards, add sources on the technologies of your
% project in alphabetical order, following the example:
% \bibitem{fastapi}
% FastAPI. [Online]. Available: \url{https://fastapi.tiangolo.com} (accessed: Mar.~14, 2026).
\renewcommand{\refname}{}
\begin{thebibliography}{99}
\vspace{-1cm}

\bibitem{gost19101}
\emph{GOST 19.101-77. Types of Programs and Program Documents}. Unified System for Program Documentation.~--- Moscow, Russia: IPK Izdatelstvo Standartov, 2001.

\bibitem{gost19103}
\emph{GOST 19.103-77. Designation of Programs and Program Documents}. Unified System for Program Documentation.~--- Moscow, Russia: IPK Izdatelstvo Standartov, 2001.

\bibitem{gost19104}
\emph{GOST 19.104-78. Basic Inscriptions}. Unified System for Program Documentation.~--- Moscow, Russia: IPK Izdatelstvo Standartov, 2001.

\bibitem{gost19105}
\emph{GOST 19.105-78. General Requirements for Program Documents}. Unified System for Program Documentation.~--- Moscow, Russia: IPK Izdatelstvo Standartov, 2001.

\bibitem{gost19106}
\emph{GOST 19.106-78. Requirements for Program Documents Produced by Printed Means}. Unified System for Program Documentation.~--- Moscow, Russia: IPK Izdatelstvo Standartov, 2001.

\bibitem{gost19504}
\emph{GOST 19.504-79. Programmer's Manual. Requirements for Content and Presentation}. Unified System for Program Documentation.~--- Moscow, Russia: IPK Izdatelstvo Standartov, 2001.

\bibitem{gost19604}
\emph{GOST 19.604-78. Rules for Making Amendments to Program Documents Produced by Printed Means}. Unified System for Program Documentation.~--- Moscow, Russia: IPK Izdatelstvo Standartov, 2001.

\end{thebibliography}

\clearpage

% ============================================================
% APPENDIX 1 — EXAMPLE OF PROGRAMMER'S COMMANDS
% ============================================================
\section*{APPENDIX 1}
\addcontentsline{toc}{section}{APPENDIX 1}

\begin{center}
\textbf{EXAMPLE OF PROGRAMMER'S COMMANDS}
\end{center}

% Hint: provide typical programmer's commands (launching infrastructure,
% launching services, checks, stopping the environment).
\begin{verbatim}
# Launch the development infrastructure
<infrastructure launch command>

# Launch all services
<service launch command>

# Run checks
<check command>

# Stop the services
<stop command>
\end{verbatim}

\clearpage


% ============================================================
% ЛИСТ РЕГИСТРАЦИИ ИЗМЕНЕНИЙ
% ============================================================
\input{../common/change_log}

\end{document}
